%======================================================================
\chapter{Simple closed curves and their interiors}\label{ch:curves}
%======================================================================

This chapter follows F.A.\ \S\S6--7 (Theorems~20--43). It develops, from Moore's axioms, the plane
topology of simple closed curves: every simple closed curve has exactly one interior; a cross-cut
divides an interior into two interiors (the $\theta$-curve theorems); every point of a curve is
accessible from both sides; there are closed curves about every point inside every region; and two
interiors that meet can be merged into one (Theorems~41 and~42). These are the tools for the double
rulings of Chapter~\ref{ch:rulings}.

Throughout, $S$ is a space satisfying Axioms~1--8, with the topology of the bridge
(Proposition~\ref{prop:bridge}): $S$ is metrizable, the regions form a basis of connected open sets
with compact closures, and Moore's arcs and simple closed curves are the topological ones. ``Closed
curve'' means simple closed curve. For an arc we write $AXB$ (an arc from $A$ to $B$ through $X$),
$\oarc{AXB}$ for the open arc, and $AXBCA$ for the closed curve formed by two arcs $AXB$ and $ACB$ with
only $A,B$ in common.

Many arguments of F.A.\ are long case analyses with figures. We give complete proofs of the results
that are used later and whose proofs are short, and proofs in outline of the others, indicating the
construction and the key step; the formal proofs are complete.

\section{The interior of a closed curve}

\begin{theorem}[F.A.\ Theorems~20 and~21]\label{thm:fa20-21}
\leavevmode
\begin{enumerate}[label=(\alph*)]
\item \emph{(Theorem~20.)} Every boundary point of a region $R$ is a limit point of its exterior
$S\setminus\cl R$. Hence $\fr(S\setminus\cl R)=\fr R$.
\item \emph{(Theorem~21.)} If $K,R$ are regions and $\fr R\subseteq\cl K$, then $R\subseteq K$.
\end{enumerate}
\end{theorem}

\begin{proof}
(a) Let $P\in\fr R$ and let $\bar R$ be a region about $P$. Axiom~7 gives regions $L,\bar L\subseteq\bar
R$ with $L$ in the exterior of $R$, $P\in\bar L$, and every boundary point of $R$ in $\bar L$ a boundary
point of $L$. So $P\in\fr L$, and the region $\bar L\ni P$ contains a point of $L$, which lies in the
exterior of $R$ and in $\bar R$. As $\bar R$ was arbitrary, $P$ is a limit point of the exterior.

(b) The exterior $E_K=S\setminus\cl K$ is connected (Axiom~3) and does not lie in the compact set $\cl R$
(it is unbounded, Theorem~\ref{thm:fa-elementary}(18)). It misses $\fr R\subseteq\cl K$, so by
Lemma~\ref{lem:conn-frontier} it lies in $S\setminus\cl R$. Hence $R\subseteq\cl K$. If a point of $R$ were
on $\fr K$, it would be a limit point of $E_K$ by (a); but $R$ is open and misses $E_K$. So
$R\subseteq K$.
\end{proof}

\formal{\lean{HPC.Moore.fa_theorem20}, \lean{HPC.Moore.fa_theorem21}.}

\begin{corollary}[The interior of a closed curve]\label{cor:interior}
Every closed curve $J$ is the frontier of exactly one region, its \emph{interior} $I(J)$. The
\emph{exterior} $E(J)=S\setminus\cl{I(J)}$ is open and connected, $S\setminus J=I(J)\cup E(J)$ (disjoint), and
$\fr I(J)=\fr E(J)=J$. If $J'$ is a closed curve with $J'\subseteq\cl{I(J)}$, then $I(J')\subseteq I(J)$.
\end{corollary}

\begin{proof}
Existence is Axiom~8; if $R,R'$ are regions with frontier $J$, Theorem~\ref{thm:fa20-21}(b) gives
$R\subseteq R'$ and $R'\subseteq R$. $I(J)$ is connected (Axiom~2), $E(J)$ is connected (Axiom~3), and the
rest follows from Theorem~\ref{thm:fa20-21}.
\end{proof}

\formal{\lean{HPC.Moore.existsUnique_interior}, \lean{HPC.Moore.jordan}, \lean{HPC.Moore.Space.IsInteriorOf}.}

``Within $J$'' means in $I(J)$, ``without $J$'' means in $E(J)$. In F.A.\ the interior of a closed curve
is used from Theorem~23 on but never defined; Corollary~\ref{cor:interior} supplies the definition.

\paragraph{The topological setting.} From here on the axioms enter only through the following facts, and
the formal proofs of Theorems~24--34 are carried out in an arbitrary Hausdorff, locally connected space
with a family of open sets (``regions'') satisfying them:
\begin{enumerate}[label=(T\arabic*)]
\item two distinct points of a connected open set are joined by an arc in it (Theorem~\ref{thm:fa15});
\item if $K,R$ are regions and $\fr R\subseteq\cl K$ then $R\subseteq K$ (Theorem~\ref{thm:fa20-21});
\item every closed curve is the frontier of a region (Axiom~8);
\item regions and their exteriors are connected, closures of regions are compact, $S$ is not compact.
\end{enumerate}

\formal{\lean{HPC.Moore.Gen.ArcConn}, \lean{HPC.Moore.Gen.Th21Prop}, \lean{HPC.Moore.Gen.Ax8Prop},
\lean{HPC.Moore.Gen.TInt}.}

Two constructions recur. \emph{Accessibility:} if $W$ is a connected open set, $F\in W$, $Q\in\fr W$ and
$N$ a neighbourhood of $Q$, there is an arc $FH$ with $H\in\fr W\cap N$ and $FH\setminus\set H\subseteq W$ (join
$F$ by an arc in $W$ to a point of $W$ in a small connected neighbourhood of $Q$, continue in that
neighbourhood towards $Q$, and stop at the first point off $W$). \emph{Crossing:} two such arcs from a
common point $F$, glued, give an arc $HK$ through $W$ with ends on $\fr W$ near two prescribed boundary
points.

\formal{\lean{HPC.Moore.Gen.exists_access}, \lean{HPC.Moore.Gen.exists_cross}.}

\begin{theorem}[F.A.\ Theorem~23]\label{thm:fa23}
No nonempty region lies in an arc.
\end{theorem}

\begin{proof}[Proof (outline)]
Suppose the region $R$ lies in the arc $AB$. Take $O\in R$, $O\ne A,B$ (regions are infinite), a region $K$
about $O$ whose closure misses $A$ and $B$, and an arc $AXB$ in the exterior of $K$
(Corollary~\ref{cor:fa16}) from $A$ to $B$. Cutting at a last and a first point, one finds points $A_1$, $B_1$
on $AB$ on either side of $O$ and a sub-arc of $AXB$ from $A_1$ to $B_1$ meeting $AB$ only in $A_1,B_1$; with
the interval $A_1OB_1$ of $AB$ it forms a closed curve $J$. $O$ is a boundary point of $I(J)$, so a small region
about $O$, inside $R$ and missing the rest of $J$, meets $I(J)$; but it lies in $AB\subseteq J$, which misses
$I(J)$.
\end{proof}

\formal{\lean{HPC.Moore.fa_theorem23}.}

\section{Theta-curves: Theorems 24--27}

In this section $J=ACB\cup ADB$ is a closed curve with interior $R$ ($A,B$ separate $C,D$ on $J$), and
$AXB$ is an arc from $A$ to $B$ with $\oarc{AXB}$ disjoint from $J$. Put $R_1=I(AXBCA)$, $R_2=I(AXBDA)$.

\begin{theorem}[F.A.\ Theorem~24]\label{thm:fa24}
If $\oarc{AXB}\subseteq R$, then (1) $R_1\subseteq R$; (2) $\oarc{ADB}$ lies without $AXBCA$; (3)
$R_1\cap R_2=\emptyset$.
\end{theorem}

\begin{proof}
(1) $\fr R_1=AXBCA\subseteq\cl R$; apply (T2). (2) $\oarc{ADB}$ misses $\cl{R_1}=R_1\cup AXBCA$: it lies on $J=\fr R$, so it misses the open set $R\supseteq R_1$;
and it misses $AXBCA$, since $\oarc{AXB}\subseteq R$ misses $J$ and $ACB\cap ADB=\set{A,B}$. So it lies in $E(AXBCA)$.
(3) If $R_1\cap R_2\ne\emptyset$: $R_2\subseteq R$ by (1) for $R_2$, so $R_2$ misses $J\supseteq ACB$,
and it misses $\oarc{AXB}\subseteq\fr R_2$; so the connected set $R_2$ misses $\fr R_1=AXBCA$ and meets $R_1$,
whence $R_2\subseteq R_1$ (Lemma~\ref{lem:conn-frontier}). Then $D\in\fr R_2\subseteq\cl{R_1}$, against (2).
\end{proof}

\formal{\lean{HPC.Moore.fa_theorem24}, \lean{HPC.Moore.Gen.th24_sub}, \lean{HPC.Moore.Gen.th24_ext},
\lean{HPC.Moore.Gen.th24_disj}.}

\begin{theorem}[F.A.\ Theorem~25]\label{thm:fa25}
If $\oarc{AXB}\subseteq R$, then $R=\oarc{AXB}\cup R_1\cup R_2$.
\end{theorem}

\begin{proof}
``$\supseteq$'' is Theorem~\ref{thm:fa24}(1). The points of $R$ on $\fr R_1=AXBCA$ are those of $\oarc{AXB}$, so
$R\cap\cl{R_1}=R_1\cup\oarc{AXB}$, and likewise for $R_2$. Hence, with $Y=R\setminus(\oarc{AXB}\cup\cl{R_1}\cup\cl
{R_2})$, the set $R\setminus\oarc{AXB}$ is the union of the three disjoint open sets $R_1,R_2,Y$. Suppose
$Y\ne\emptyset$, and take $E\in Y$.

\emph{Step 1.} An arc in $R$ from $E$ to a point of $R_1$ (T1) must meet $\oarc{AXB}$; let $O$ be its first point on
the closed set $AXB$ ($A,B\notin R$). The sub-arc $EO$ minus $O$ is connected, lies in $R\setminus\oarc{AXB}$ and
contains $E$, so it lies in $Y$.

\emph{Step 2.} $E$ and the points of $\oarc{ADB}$ lie in the connected open set $E(AXBCA)$ (by
Theorem~\ref{thm:fa24}(2), and $E\notin\cl{R_1}$). An arc in it from $E$ to a point of $\oarc{ADB}$ has a first
point $L$ on $J$. The part before $L$ is connected, misses $J=\fr R$ and starts in $R$, so lies in $R$; it misses
$\cl{R_1}\supseteq\oarc{AXB}$; so it lies in $R_2\cup Y$, hence in $Y$. And $L\in J\setminus\cl{R_1}=\oarc{ADB}$.

\emph{Step 3.} Choose points $\bar H,H,O,K,\bar K$ in this order on $\oarc{AXB}$ and, by the crossing construction, arcs
$HP_1K$ with open part in $R_1$ and $\bar HP_2\bar K$ with open part in $R_2$ (the ends can be prescribed near
given points of $\oarc{AXB}$ because these points are on $\fr R_1$ and $\fr R_2$; one aims at points of $\oarc{AXB}$
on either side of $O$, which is why the order can be arranged). Then
$\bar J=\bar HP_2\bar K\cup\bar KK\cup KP_1H\cup H\bar H$ (with $\bar KK$, $H\bar H$ sub-arcs of $AXB$) is a closed curve
lying in $R$, so its interior $\bar R$ lies in $R$ by (T2).

\emph{Step 4.} $O\notin\bar J$, and $O$ is joined to $L$ by $OE\cup EL$, which misses $\bar J$ (its points other than $O,L$
lie in $Y$, and $L\notin R$). As $L\notin\cl{\bar R}\subseteq R$, $O$ lies in $E(\bar J)$ (Lemma~\ref{lem:conn-frontier}).
The open sub-arc of $AXB$ between $H$ and $K$ misses $\bar J$ and contains $O$, so it lies in $E(\bar J)$; the open
sub-arcs from $A$ to $\bar H$ and from $\bar K$ to $B$ miss $\bar J$ and have the ends $A,B\notin\cl{\bar R}$, so they lie
in $E(\bar J)$ too. Hence $\bar R$ misses $\oarc{AXB}$ and $\bar R\subseteq R_1\cup R_2\cup Y$. But $\bar R$ is connected and
meets $R_1$ (near $P_1\in\fr\bar R\cap R_1$) and $R_2$ (near $P_2$): a contradiction.
\end{proof}

\formal{\lean{HPC.Moore.fa_theorem25}, \lean{HPC.Moore.Gen.th25}.}

\begin{theorem}[F.A.\ Theorem~26]\label{thm:fa26}
If $ACB$, $ADB$, $AFB$ are three arcs, no two with a common point other than $A,B$, then two of the interiors of
$ACBFA$, $ACBDA$, $ADBFA$ meet.
\end{theorem}

\begin{proof}[Proof (outline)]
Suppose the three interiors $R=I(ACBFA)$, $R_1=I(ACBDA)$, $R_2=I(ADBFA)$ are pairwise disjoint. Then each open arc
lies without the curve formed by the other two (for $\oarc{ADB}$: were it within $ACBFA$, Theorem~\ref{thm:fa24} would put
$R_1$ inside $R$). The set $Y$ of points without all three curves is nonempty (the closures of the three interiors are
compact, $S$ is not). From a point $P\in Y$, an arc without $ACBFA$ reaches $\oarc{ADB}$ at a first point $O$, through $Y$.
Around $O$ one builds, by crossing arcs through $R_1$ and $R_2$, a closed curve $\bar J$ on $R_1\cup R_2\cup\oarc{ADB}$ as in
the proof of Theorem~\ref{thm:fa25}. If $I(\bar J)$ met $ACB\cup AFB$ it would contain it and then all of $Y$, which is
unbounded; if it met $Y$ it would meet $ACB$ through arcs in $Y$; so $I(\bar J)\subseteq R\cup R_1\cup R_2$, while it meets
$R_1$ and $R_2$ and is connected. Contradiction. (This simplifies Moore's three cases: his Case~I is reduced to Case~II and
his auxiliary curves are not needed.)
\end{proof}

\formal{\lean{HPC.Moore.fa_theorem26}, \lean{HPC.Moore.Gen.th26}.}

\begin{theorem}[F.A.\ Theorem~27]\label{thm:fa27}
If $\oarc{AXB}$ lies without $J$, then (1) $D$ is without $AXBCA$ or $C$ is without $AXBDA$; (2) if $D$ is without
$AXBCA$, then $C$ is within $AXBDA$ and $I(AXBDA)=\oarc{ACB}\cup R\cup R_1$.
\end{theorem}

\begin{proof}
(1) If $D$ is not without $AXBCA$, then $\oarc{ADB}$ (connected, missing $AXBCA$) is within $AXBCA$, and
Theorem~\ref{thm:fa24}(2), applied to the curve $AXBCA$ and the cross-cut $\oarc{ADB}$, puts $C$ without $AXBDA$.
(2) $R$ and $R_1$ are disjoint ($R_1$ is connected, misses $J$, and is not inside $R$, as its boundary point $X$ is
without $J$). If $C$ were also without $AXBDA$, the three interiors $R,R_1,R_2$ would be pairwise disjoint
(the cases are symmetric), against Theorem~\ref{thm:fa26}. So $C$ is within $AXBDA$; then $\oarc{ACB}$ is a
cross-cut of $AXBDA$ and Theorem~\ref{thm:fa25} gives $I(AXBDA)=\oarc{ACB}\cup R\cup R_1$.
\end{proof}

\formal{\lean{HPC.Moore.fa_theorem27}, \lean{HPC.Moore.Gen.th27}.}

\begin{corollary}[Splitting an interior]\label{cor:split-interior}
If $J$ is a closed curve, $AXB$ an arc with $A,B\in J$ and $\oarc{AXB}\subseteq I(J)$, and $C_1,C_2$ the two arcs of $J$
from $A$ to $B$, then $I(J)\setminus AXB$ is the disjoint union of $I(AXB\cup C_1)$ and $I(AXB\cup C_2)$.
\end{corollary}

\formal{\lean{HPC.Moore.split_interior}, \lean{HPC.Moore.split_interior_of}.}

\section{Arcs near a curve: Theorems 28--30}

\begin{theorem}[F.A.\ Theorem~28]\label{thm:fa28}
Let $O$ be a point of the closed curve $J$, $\bar R$ a region about $O$, and $M$ either $I(J)$ or $E(J)$. There is an
arc $AXB$ with $A,B\in J$, $\oarc{AXB}\subseteq M\cap\bar R$, such that of the two arcs of $J$ from $A$ to $B$ the one
containing $O$ (in its interior) lies in $\bar R$.
\end{theorem}

\begin{proof}[Proof (outline)]
Axiom~6 (for $M=I(J)$) or Axiom~7 (for $M=E(J)$) gives regions $L\subseteq M\cap\bar R$ and $\bar L\subseteq\bar R$ about $O$ such
that every point of $J$ in $\bar L$ is a boundary point of $L$. Take an arc $A_1OB_1$ of $J$ in $\bar L$, points $A_2,B_2$
with $A_1A_2OB_2B_1$ in order, small regions $R_A\ni A_2$, $R_B\ni B_2$ meeting $J$ only in $A_1O$ and $OB_1$, points
$C\in L\cap R_A$, $D\in L\cap R_B$, an arc from $C$ to $D$ in $L$ and arcs from $C$ to $A_2$ in $R_A$, from $D$ to $B_2$ in
$R_B$. The arc $AXB$ is extracted from their union; Moore's last step needs care, since the arc $CA_2$ may leave $M$
before reaching $J$: one takes for $A$ the \emph{first point of $CA_2$ outside $M$}, which lies on $\fr M=J$, hence on the
open arc $A_1O$; likewise $B$.
\end{proof}

\formal{\lean{HPC.Moore.fa_theorem28}.}

Theorem~29 (if $A,B$ separate $C,D$ on $J$, and $\oarc{AXB}$, $\oarc{CYD}$ are both within or both without $J$, they meet)
follows from Theorems~\ref{thm:fa24}, \ref{thm:fa25} and~\ref{thm:fa27}. In the case ``within'', if the arcs did not
meet, the connected set $\oarc{CYD}$ would lie in $R_1\cup R_2$ (Theorem~\ref{thm:fa25}), hence in one of them, say $R_1$; but its
limit point $D$ lies without $AXBCA$ (Theorem~\ref{thm:fa24}(2)), and symmetrically $C$ lies without $AXBDA$. The case ``without'' uses
Theorem~\ref{thm:fa27} in the same way. Theorem~30 (the interiors of $AXBDA$ and $CYDBC$ meet) is not proved by Moore; the formal proof
examines the four combinations of the alternatives of Theorem~\ref{thm:fa27}.

\formal{\lean{HPC.Moore.fa_theorem29}, \lean{HPC.Moore.fa_theorem30}.}

\section{The two sides of an arc and accessibility: Theorems 31--34}

\begin{theorem}[F.A.\ Theorem~31; the pocket construction]\label{thm:fa31}
Let $f\colon I\to S$ parametrize an arc, $0<b<1$, and $R$ a region about $f(b)$. There are closed curves $J_1=F_1\cup P_1$
and $J_2=F_2\cup P_2$, where $P_1\supseteq P_2$ are sub-arcs of the arc around $f(b)$ (in its interior) and $F_1,F_2$ are
arcs meeting the given arc only in their ends, such that $I(J_1)$ and $I(J_2)$ are disjoint, miss the arc, and together
with the open arc $\oarc{P_2}$ form a region about $f(b)$ inside $R$.
\end{theorem}

\begin{proof}[Proof (outline)]
A closed curve $J_0$ through $f(b)$ containing a sub-arc of the given arc is built as in Theorem~\ref{thm:fa23}. Theorem~28
inside $J_0$ at $f(b)$ gives an arc $F_1$ within $J_0$, closing a curve $J_1$; Theorem~28 again, for $J_1$ from outside,
gives $F_2$ without $J_1$, closing a curve $J_2$. By Theorem~\ref{thm:fa27}, the interior of the curve formed by $F_1$, $F_2$
and the two short pieces of the arc is $I(J_1)\cup I(J_2)\cup\oarc{P_2}$, and Theorem~\ref{thm:fa24}(3) makes $I(J_1)$ and
$I(J_2)$ disjoint.
\end{proof}

\formal{\lean{HPC.Moore.fa_theorem31}, \lean{HPC.Moore.th31_pocket}.}

\begin{corollary}[The two sides]\label{cor:sides}
If $f(b)$ is an interior point of an arc $\alpha$ and $W$ an open set about $f(b)$, there are an open set $N$ with $f(b)\in
N\subseteq W$ and two connected open sets $I_1,I_2\subseteq W\setminus\alpha$ with $N\setminus\alpha\subseteq I_1\cup I_2$, each of
which has every point of some open sub-arc of $\alpha$ about $f(b)$ as a limit point.
\end{corollary}

\formal{\lean{HPC.Moore.sides}.}

\begin{theorem}[Propagation and accessibility; F.A.\ Theorem~32]\label{thm:fa32}
Let $f$ parametrize an arc $\alpha=AB$ ($A=f(0)$, $B=f(1)$), $O$ an open set, and $Q$ the component of $O\setminus\alpha$
containing a point $x_0$.
\begin{enumerate}[label=(\alph*)]
\item \emph{(Propagation.)} If $T\subseteq(0,1)$ is an interval with $f(T)\subseteq O$ and some point of $f(T)$ is a limit point
of $Q$, then every point of $f(T)$ is a limit point of $Q$.
\item \emph{(Accessibility.)} If $f(u)\in O$ for all $u\in(t_0,1)$ and some such $f(u)$ is a limit point of $Q$, then $x_0$ is
joined to $B$ by an arc lying, except for $B$, in $Q$.
\end{enumerate}
\end{theorem}

\begin{proof}
(a) The set of $u\in T$ with $f(u)\in\cl Q$ is closed in $T$ and nonempty. It is open: if $f(u)\in\cl Q$, take $N,I_1,I_2$
as in Corollary~\ref{cor:sides} with $W\subseteq O$. $Q$ meets $N\setminus\alpha$, hence $I_1$ or $I_2$, say $I_1$; as $I_1$ is
connected and lies in $O\setminus\alpha$, $I_1\subseteq Q$, and every point of an open sub-arc about $f(u)$ is in $\cl{I_1}\subseteq\cl
Q$. As $T$ is connected, the set is all of $T$.

(b) (Outline.) By (a), $Q$ has every point of $f((t_0,1))$ as a limit point. Choose open sets $O_k\subseteq O$, $O_k$ within
distance $1/k$ of $B$ and containing $f((t_k,1))$ for suitable $t_k\to1$; by (a) the component $Q_k$ of $O_k\setminus\alpha$ adjacent to
$f((t_k,1))$ lies in $Q$ and, by connectedness, in $Q_{k-1}$. Pick $x_k\in Q_k$ and arcs $\beta_k$ in $Q_k$ from $x_k$ to
$x_{k+1}$ (T1), arranged so that non-consecutive ones are disjoint. The $\beta_k$ converge to $B$, and
Theorem~\ref{thm:limit-arc} (with trimming) gives an arc from $x_0$ to $B$ in $Q\cup\set B$.
\end{proof}

\formal{\lean{HPC.Moore.side_propagate}, \lean{HPC.Moore.accessible}, \lean{HPC.Moore.fa_theorem32}.}

Moore's own proof of Theorem~32 covers the arc by chains of the closed curves of Theorem~31 and uses a least-upper-bound
argument with two classes of points; the connectedness argument (a) replaces it.

\begin{theorem}[F.A.\ Theorems~33 and~34]\label{thm:fa33-34}
\leavevmode
\begin{enumerate}[label=(\alph*)]
\item If $A\in\fr R$ for a region $R$, $B\in R$, and $AB$ is an arc with $AB\setminus\set A\subseteq R$, then
$R\setminus(AB\setminus\set A)$ is connected.
\item If $AB$ is an arc lying in a region $R$, then $R\setminus AB$ is connected. In particular this holds for $R$ the
interior of a closed curve.
\end{enumerate}
\end{theorem}

\begin{proof}
We use the following consequence of Theorem~\ref{thm:fa32} (\emph{reaching an end}): if $O$ is open, $\alpha=AB$ an arc
with $\alpha\setminus\set A\subseteq O$, and $P\in O\setminus\alpha$ is joined in $O$ by an arc to a point of $\alpha$ other
than $A$, then $P$ is joined to $B$ by an arc in $O$ meeting $\alpha$ only in $B$. Indeed, the first point $X$ of that arc on
$\alpha$ is a limit point of the component $Q$ of $P$ in $O\setminus\alpha$; if $X=B$ we are done, otherwise propagate along
$\alpha$ from $X$ towards $B$ and apply accessibility.

(a) (Kline's proof.) Let $C,D\in R\setminus(AB\setminus\set A)$. An arc from $C$ to $B$ in $R$ meets $AB$, not at
$A\notin R$; so $C$ is joined to $B$ by an arc $\gamma$ meeting $AB$ only in $B$. Then $\Gamma=\gamma\cup AB$ is an arc from
$C$ to $A$ with $\Gamma\setminus\set A\subseteq R$. An arc from $D$ to $C$ in $R$ meets $\Gamma$ not at $A$; reaching the end $C$
of $\Gamma$, $D$ is joined to $C$ by an arc in $R$ meeting $\Gamma$ only in $C$. This arc misses $AB$, so $C$ and $D$ are in the
same component of $R\setminus(AB\setminus\set A)$.

(b) The same argument, with both ends of $AB$ in $R$: reaching an end gives an arc from $C$ to an end of $AB$, say $B$,
meeting $AB$ only there, and $\Gamma=\gamma\cup AB$; then $D$ reaches an end of $\Gamma$. If it reaches $C$ we are done. If it
reaches $A$, then $D$ and $C$ are the ends of an arc $D\dots A\dots B\dots C$ in $R$ containing $AB$, and the two sides of this
arc near an interior point (Corollary~\ref{cor:sides}) together with accessibility join $D$ to $C$ in $R\setminus AB$.
\end{proof}

\formal{\lean{HPC.Moore.reach_end}, \lean{HPC.Moore.joinZ}, \lean{HPC.Moore.fa_theorem33},
\lean{HPC.Moore.fa_theorem34_region}, \lean{HPC.Moore.fa_theorem34}.}

Moore gives two proofs of Theorem~34. The first (for interiors of closed curves) is not needed; the second, ``a
modification of a proof due to Mr.~J.~R.~Kline'', is the one formalized. Kline's argument does not cover the case in
which the arc from $D$ first meets $\Gamma$ at its end $A$; that case is the last sentence of the proof.

\section{Curves about points: Theorems 35, 36, 39 and 40}

\begin{theorem}[F.A.\ Theorems~39 and~40; accessibility of a closed curve]\label{thm:fa39-40}
Let $J$ be a closed curve and $M$ either $I(J)$ or $E(J)$.
\begin{enumerate}[label=(\alph*)]
\item Every point $O\in J$ is joined to every point $P\in M$ by an arc lying in $M\cup\set O$. In particular every point not on
$J$ is joined to $O$ by an arc meeting $J$ only in $O$.
\item Any two points $A\ne B$ of $J$ are joined by an arc whose open part lies in $I(J)$ and by one whose open part lies in
$E(J)$.
\end{enumerate}
\end{theorem}

\begin{proof}
(a) Let $f$ parametrize an arc $\alpha$ of $J$ with $f(1)=O$. The set $W=(S\setminus J)\cup f((0,1))$ is open (its
complement is the compact set $J\setminus f((0,1))$), it contains $f(u)$ for $0<u<1$, and $W\setminus\alpha=S\setminus J$. So the
component of $P$ in $W\setminus\alpha$ is the component of $P$ in $S\setminus J$, which is $M$ (Corollary~\ref{cor:interior}).
Since $\fr M=J$, the points $f(u)$ are limit points of $M$. Theorem~\ref{thm:fa32}(b) joins $P$ to $O$ by an arc lying, except for
$O$, in $M$.
(b) Take $P\in M$ and arcs from $P$ to $A$ and to $B$ as in (a); their union contains an arc from $A$ to $B$
(Lemma~\ref{lem:arc-ops}(c)), whose open part lies in $M$.
\end{proof}

\formal{\lean{HPC.Moore.access}, \lean{HPC.Moore.fa_theorem39}, \lean{HPC.Moore.fa_theorem40}.}

Moore proves Theorem~39 with a sequence of closed curves converging to $O$ and a limit arc; accessibility
(Theorem~\ref{thm:fa32}) does the same work.

\begin{theorem}[F.A.\ Theorems~35 and~36]\label{thm:fa35-36}
\leavevmode
\begin{enumerate}[label=(\alph*)]
\item Every nonempty region contains a closed curve.
\item If $O$ is a point of a region $R$, there is a closed curve $J\subseteq R$ with $O\in I(J)\subseteq R$.
\end{enumerate}
\end{theorem}

\begin{proof}[Proof (outline)]
(a) Take an arc in $R$ (Corollary~\ref{cor:fa16}) and, about an interior point of it, the curve $J_1$ of the pocket
construction (Theorem~\ref{thm:fa31}), which lies in $R$.

(b) Moore's proof, with one correction. Take a closed curve $J$ in a small region about a point $X\ne O$ of $R$, so that $O$ is
without $J$ and $I(J)\subseteq R$ (by (a) and (T2)); an arc $OA$ meeting $J$ only in $A$; a boundary point $C$ of $R$
(Theorem~\ref{thm:fa-elementary}(19)) and an arc $BC$ from $J$ to $C$ with $BC\setminus\set C\subseteq R$, missing $OA$ and meeting $J\cup I(J)$
only in $B$ (by Theorem~\ref{thm:fa33-34}(a) in $R\setminus OA$ and first points). Moore takes $BC$ inside a region about $C$ that need
not lie in $R$; the arc is cut at its first point outside $R$, which is again a boundary point of $R$, and that point is used as $C$.
Let $AF_1B$, $AH_1B$ be the two arcs of $J$ from $A$ to $B$ and $FZH$ an arc from $AF_1B$ to $AH_1B$, missing $OA\cup BC$, with open part
without $J$ and inside $R$ (Theorems~\ref{thm:fa20-21}, \ref{thm:fa33-34}). Then $B$ is without $ZHAFZ$ (the interior of that curve lies
in $R$, and the arc $BC$ reaches $C\notin R$ without meeting it), so by Theorem~\ref{thm:fa27} $A$ is within $ZHBFZ$, and so is $O$, since
$OA$ does not meet that curve. Finally $I(ZHBFZ)\subseteq R$ by (T2), as the curve lies in $R$.
\end{proof}

\formal{\lean{HPC.Moore.fa_theorem35}, \lean{HPC.Moore.fa_theorem36},
\lean{HPC.Moore.fa_theorem36_interior_subset}.}

\begin{remark}
As printed, Theorem~35 (``every region contains a closed curve'') fails for the empty region
(\lean{HPC.Moore.not_fa_theorem35_printed}); see Remark~\ref{rem:empty-region}. The statement that the interior lies in $R$ is
not part of the printed Theorem~36, but it is what Moore uses in 1925 (Theorem~19), and it follows from (T2).
\end{remark}

\section{Merging interiors: Theorems 37--43}

\begin{definition}\label{def:chord}
Let $J_1,J_2$ be closed curves. A \emph{chord} of $J_1$ over a closed set $K$ is an arc of $J_1$ with its ends on $K$ and its open
arc off $K$.
\end{definition}

Every point of $J_1\setminus K$ lies on exactly one chord over $K$ (cut $J_1$ at the nearest points of $K$ on either side), and two
chords whose open arcs meet are equal.

\begin{lemma}[Chords form a null family; Moore 1919, Theorem~A]\label{lem:chords-null}
Arcs of a closed curve $L$ with pairwise disjoint open arcs form a null family: for every $\varepsilon>0$ only finitely many of
them have a point at distance $\ge\varepsilon$ from their first end.
\end{lemma}

\begin{proof}
Parametrize $L$ by the circle $\R/\Z$. Each arc of $L$ is the image of a closed interval of the circle, and the open arcs being
disjoint, so are the open intervals; so for every $\delta>0$ only finitely many intervals have length $\ge\delta$. By uniform
continuity of the parametrization, an arc whose interval has length $<\delta$ has diameter $<\varepsilon$ for suitable $\delta$.
\end{proof}

\formal{\lean{HPC.Moore.chords_null}, \lean{HPC.Moore.m19_theoremA}.}

Moore's Theorem~A (1919) is the statement that a closed curve $L$ carries no infinite family of arcs, no two with a common
segment, each meeting two given disjoint closed sets; the lemma above is the form used here.

Theorems~37 and~38 of F.A.\ produce, for a point $O$ of $J$ not on another closed curve $C$, the innermost (resp.\ outermost) arc of
$C$ around $O$ with open part within (resp.\ without) $J$. The formal proofs replace Moore's limit-point arguments by the
observation that the chords of $C$ in question have only finitely many pairs of ends: distinct chords have disjoint open
intervals in a parametrization of $C$, and all have length bounded below, since their two ends lie in disjoint compact sets.

\formal{\lean{HPC.Moore.exists_maximal_chord}, \lean{HPC.Moore.fa_theorem37}, \lean{HPC.Moore.fa_theorem38}.}

\begin{theorem}[F.A.\ Theorem~41]\label{thm:fa41}
If the interiors $I_1,I_2$ of the closed curves $J_1,J_2$ meet, there is a closed curve $J\subseteq J_1\cup J_2$ with $I_1\cup
I_2\subseteq I(J)$.
\end{theorem}

\begin{proof}
\emph{Case I: $J_1\cap J_2$ has at most one point.} Then one interior contains the other, and $J$ is the bigger curve. Indeed,
if $J_2$ misses $I_1$, then $I_1$ (connected, meeting $I_2$, missing $J_2=\fr I_2$) lies in $I_2$. Otherwise $J_2$ has a point $x$
within $J_1$, and no point $y$ without $J_1$: the two arcs of $J_2$ from $x$ to $y$ would each meet $J_1$ at a point other than
$x,y$, giving two common points. So $J_2\subseteq\cl{I_1}$ and $I_2\subseteq I_1$ by (T2).

\emph{Case II, first sub-case: every point of $J_1$ is within or on $J_2$.} Then $I_1\subseteq I_2$ by (T2), and $J=J_2$;
symmetrically if $J_2\subseteq\cl{I_1}$.

\emph{Case II proper} (outline). Let $V$ be a component of $S\setminus(\cl{I_1}\cup\cl{I_2})$ with noncompact closure. Its frontier lies
on $J_1\cup J_2$, and $V$ misses the interior of every closed curve that misses $V$. There is a point $O$ of $J_2$ without $J_1$ which is
a limit point of $V$, and a second one $O'$. Split $J_2$ at $O,O'$ into two arcs; on each, replace every sub-arc cut off by a chord of $J_1$
over $J_2$ lying in $\cl V$ by that chord. The replaced sub-arcs have disjoint parameter intervals (no point inside such a sub-arc is a limit
point of $V$, by Theorem~\ref{thm:fa27}), the chords have disjoint open arcs and form a null family (Lemma~\ref{lem:chords-null}), so the
result is an arc from $O$ to $O'$ (Theorem~\ref{thm:arc-replace}). The two arcs so obtained meet only in $O,O'$, and their union $J$ is a
closed curve on $J_1\cup J_2$ containing the frontier of $V$. Then $V$ is the exterior of $J$, and $I_1\cup I_2$, which misses $J$ and $V$, lies
in $I(J)$.
\end{proof}

\formal{\lean{HPC.Moore.fa_theorem41_caseI_nested}, \lean{HPC.Moore.fa_theorem41_caseI},
\lean{HPC.Moore.fa_theorem41_of_subset}, \lean{HPC.Moore.fa_theorem41_caseII}, \lean{HPC.Moore.Gen.arc_replace},
\lean{HPC.Moore.fa_theorem41}.}

Moore ends the proof of Case~II with ``It may be proved that $J$ is a closed curve satisfying the requirements of Theorem~41''; the
replacement of a null family of arcs (Theorem~\ref{thm:arc-replace}) is what this sentence needs.

\begin{theorem}[F.A.\ Theorem~42]\label{thm:fa42}
Let $\mathcal G$ be a finite nonempty set of closed curves the union of whose interiors is connected. There is a closed curve lying on
the curves of $\mathcal G$ whose interior contains all their interiors.
\end{theorem}

\begin{proof}
Induction on the number of curves not yet absorbed. Suppose $\mathcal H\subsetneq\mathcal G$ is nonempty and $J_{\mathcal H}$ is a closed
curve on the curves of $\mathcal H$ whose interior contains theirs. If no curve of $\mathcal G\setminus\mathcal H$ had an interior meeting
$U=\bigcup_{\mathcal H}I(J)$, then $U$ and $V=\bigcup_{\mathcal G\setminus\mathcal H}I(J)$ would be disjoint nonempty open sets whose union
is the union of all the interiors, against connectedness. So some curve $q\in\mathcal G\setminus\mathcal H$ has an interior meeting
$U\subseteq I(J_{\mathcal H})$. Theorem~\ref{thm:fa41}
applied to $J_{\mathcal H}$ and $q$ absorbs $q$.
\end{proof}

\formal{\lean{HPC.Moore.fa_theorem42_of_41}, \lean{HPC.Moore.fa_theorem42}.}

The connectedness of the union of the interiors is essential; it is what guarantees an order in which each new interior meets the union
of the previous ones.

\begin{theorem}[F.A.\ Theorem~43, corrected]\label{thm:fa43}
Let $J$ be a closed curve, $O\in J$, and $\mathcal G$ a finite \emph{nonempty} set of closed curves whose interiors contain $O$ and none
of whose interiors contains $I(J)$. There are closed curves $J_1,J_2$ through $O$, lying on $J$ and the curves of $\mathcal G$, with $I(J_1)$
without $J$, $I(J_2)$ within $J$, and both inside every interior of a curve of $\mathcal G$.
\end{theorem}

\begin{proof}[Proof (outline)]
Moore gives no proof. The key step is the intersection analogue of Theorem~41: if $O\in J$ lies within $J'$ and $U$ is $I(J)$ or $E(J)$
(in the second case with $I(J)\not\subseteq I(J')$), and $J\cap J'$ has two points, then the component of $U\cap I(J')$ adjacent to $O$ is
the interior of a closed curve on $J\cup J'$ through $O$. Its proof is that of Case~II of Theorem~41 with this component in place of the
unbounded one. Iterating over $\mathcal G$ gives $J_2$ (starting from $J$) and $J_1$ (starting from the exterior side).
\end{proof}

\formal{\lean{HPC.Moore.dual_component}, \lean{HPC.Moore.fa_theorem43}.}

For $\mathcal G=\emptyset$ the printed Theorem~43 is false: its hypotheses hold vacuously and its conclusion would make the interior of $J$
lie in its own exterior (\lean{HPC.Moore.not_fa_theorem43_printed}).

\section*{Exercises}
\addcontentsline{toc}{section}{Exercises}

\begin{exercise}
In $\R^2$, verify Theorem~\ref{thm:fa25} for $J$ the unit circle and $AXB$ a diameter, and Theorem~\ref{thm:fa27} for $AXB$ a semicircle
outside the unit circle.
\end{exercise}

\begin{exercise}
Show by an example in $\R^2$ that Theorem~\ref{thm:fa42} fails if the union of the interiors is not connected.
\end{exercise}

\begin{exercise}
Prove Corollary~\ref{cor:split-interior} from Theorems~\ref{thm:fa24} and~\ref{thm:fa25}.
\end{exercise}
